\section[概览与模形式的一种视角]{概览与模形式的一种视角（1 月 19 日）}
\srcpage{2}

这是一门关于几何 Langlands 纲领的课程，采取一个特定的视角，并把它与电磁对偶（electric-magnetic duality）的关系一并纳入。课程通过 Zoom 进行。

几何 Langlands 纲领位于数论与物理之间，各占一半。也许我们就是 Odysseus，正要设法航行回家，途中要穿过两重危险：一边是 Charybdis（物理，古典意义上那道大漩涡），另一边是 Scylla（数论，古典意义上那头海怪）。你大概还能把这个类比继续延伸下去，例如导出代数几何（derived algebraic geometry）就是 Calypso 的岛屿。继续延伸这个类比，就留作练习。

这门课没有特别指定的预备知识——特别是，你不需要懂物理或数论。Mackey [Mac80] 写过一篇很好的综述，讲的是对称性与调和分析之关系的一个视角，读起来大概会很有趣。本讲与下一讲，我们会谈模形式以及它与物理的一些联系；此后课程才正式开始：某种意义上，几何 Langlands 纲领是 Fourier 变换的一次巨大推广，所以我们会从 Fourier 变换讲起，讲法要便于日后推广。

\paragraph{模形式、自守形式与物理}

我们不会在数论上过于技术化。模形式与自守形式（automorphic forms）的想法，是在算术局部对称空间（arithmetic locally symmetric space）上做某种调和分析或量子力学。举例来说，上半平面 $\mathbb{H}$ 有一个模型 $\mathrm{SL}_2(\mathbb{R})/\mathrm{SO}_2$。模群 $\Gamma := \mathrm{SL}_2(\mathbb{Z})$ 作用于 $\mathbb{H}$，其基本区域（fundamental domain）记为 $\Gamma\backslash\mathbb{H}$（TODO：图）：这个基本区域非紧，沿 $y$ 轴伸向无穷，还有几个 orbifold 点，即 $\Gamma$-作用的稳定子所在之处。

更一般地，我们可以考虑带有极大紧子群 $K\subset G$ 的 Lie 群 $G$，以及格 $\Gamma\subset G$。然后考虑空间 $\Gamma\backslash G/K$，并研究其上的函数空间。你可以想象一个粒子在这个局部对称空间上运动，于是我们关心 $L^2(\Gamma\backslash G/K)$，其上有 Laplacian $\Delta$ 作用；我们可以按特征函数子空间来分解函数空间。这是模形式出现的方式之一。

这类变体多得数不清。你可以把 $K$ 一脚踢出这个故事（yeet），直接研究 $L^2(\Gamma\backslash G)$ 及其 $K$-作用。（TODO：关于单位切丛的某些话。）此外，也不必把自己限制在用函数来线性化：可以用形式，或用别的向量丛的截面，例如 $\Gamma(\Gamma\backslash\mathbb{H},\omega^{\otimes k/2})$。de Rham 说，这与 $\Gamma\backslash\mathbb{H}$ 的上同调有关，可能带扭系数。所有这些变体都是与模形式有关的对象的例子。

现在你也许在想：一旦过渡到上同调，就不再是在做量子力学了。但事实上，这正是所谓拓扑量子力学（topological quantum mechanics）的领域；例如 Witten 在关于超对称量子力学与 Morse 理论之关系的论文 [Wit82] 中讨论过这一点。

自守形式遵循类似的故事，只是 $G$ 是更一般的 Lie 群。例如，取你最喜欢的约化代数群，比如 $\mathrm{GL}_n$ 或 $\mathrm{Sp}_n$；令 $G$ 为该群的实点，$\Gamma$ 为该群的整点，$K$ 为 $G$ 的极大紧子群。用调和分析研究 $\Gamma\backslash G/K$ 上的函数空间、并把它看成量子力学，这有很长的历史。例如，若从 $\mathrm{Sp}_{2n}$ 出发，就得到 $\mathrm{Sp}_{2n}(\mathbb{Z})\backslash\mathrm{Sp}_{2n}(\mathbb{R})/U_n$。

但这里的结构比典型的量子力学设定多得多。在模形式的情形（$G=\mathrm{SL}_2(\mathbb{R})$）已经能看出这一点。即：还有一个额外的变量：我们可以把 $\mathbb{Z}$ 推广成数域的其他整数环。也就是说，给定域 $\mathbb{Q}$，把 $\mathbb{Z}$ 看作 $\mathcal{O}_{\mathbb{Q}}$，即该数域的整数环，从而得到群 $\mathrm{SL}_2(\mathcal{O}_{\mathbb{Q}})$。现在可以把 $\mathbb{Q}$ 换成任意有限扩张 $F/\mathbb{Q}$，令 $\mathcal{O}_F$ 为 $F$ 的整数环，考虑新的格 $\mathrm{SL}_2(\mathcal{O}_F)$。要把它做得完全精确，还得摆弄 $\mathrm{SL}_2(\mathbb{R})$，因为 $F$ 在无穷处可能不止一个位（place）；不过这属于我们眼下要回避的技术细节。

还有另一种改变数据的方式：固定 $F$，比如 $F=\mathbb{Q}$。然后可以改变导子（conductor）或分歧数据（ramification data）。也就是说，$\Gamma$ 在 $\mathbb{H}$ 上的基本区域有许多覆叠空间 $\Gamma'\backslash\mathbb{H}$，其中 $\Gamma'\subset\mathrm{SL}_2(\mathbb{Z})$ 是同余子群（congruence subgroup）。同余子群的一个例子是：给定 $N\in\mathbb{Z}$，子群
\[
\Gamma(N) := \left\{ \begin{pmatrix} a & b \\ c & d \end{pmatrix} \equiv \mathrm{id} \bmod N \right\}.
\]
一个变体是 $\Gamma_0(N)$，即模 $N$ 为上三角的矩阵构成的子群。

\srcpage{3}

正如数论中常见的那样，我们可以考察 $\mathcal{O}_F$ 中素元处的不同位。例如 $F=\mathbb{Q}$ 时，这意味着考察素点 $p$ 处的局部数据，其中涉及 $\mathrm{SL}_2(\mathbb{Q}_p)$。所以我们最终想造出的 Hilbert 空间依赖于所有这些数据：$G$、$\Gamma$、$K$，还可能依赖于 $F$、同余子群以及素点。

无论如何，我们会得到一个 Hilbert 空间，并可以研究 Laplacian 的谱理论。也许出人意料的是，特征空间通常不是一维的。这种「退化」来自 Hecke 算子，它是这个故事的关键部分。粗略地说，Laplacian 属于一个由两两交换的算子组成的大族；若 $p$ 是不整除 $N$ 的素数，这个族有一个成员 $T_p$，称为 Hecke 算子，它给出 $\mathbb{C}[T_p]$ 在 $L^2(\Gamma\backslash G/K)$ 上的作用。这些算子全都交换，所以把所有素数的 $\mathbb{C}[T_p]$ 张量起来，就作用在 Hilbert 空间上，并保持各特征空间。

从量子力学的视角看，这么多交换算子是很不寻常的。你可以把它想成一个可积系统（integrable system），带有大量守恒量。通常（TODO：如果我理解得没错），可积系统是混沌的反面，但这些算术系统却被当作量子混沌（quantum chaos）的好例子来研究！这是算术故事的一个特点：「算术量子混沌」的行为，比人们预想的更像量子可积系统。

在这个系统中，对模形式（或自守形式）有一族特殊的测量/状态，称为周期（periods）。一个基本的例子是：给定基本区域 $\Gamma\backslash\mathbb{H}$ 上的模函数 $f$，对它积分：
\[
\int_{i\mathbb{R}_+} f.
\]
我们会结合这些不变量来研究模函数/模形式。Hecke 用 $L$-函数造出了这类不变量的例子。

\begin{definition}
Maass 形式（Maass form）是 $L^2(\Gamma\backslash G/K)$ 上 Laplacian 的特征函数。具体地说，模形式是 $\omega^{\otimes k/2}$ 的全纯截面。
\end{definition}

模形式也可以通过考察 $\Gamma\backslash\mathbb{H}$ 的（扭）上同调得到；这就是所谓的 Eichler–Shimura 理论。

本课程的重点更多在拓扑量子力学，而不是量子力学；我们主要关心基态。在这里，模形式有点像是基态。

在已经堆成一大摞的结构之上，还要再加一块本质性的结构。有一些神秘的算子，可以让你改变群！也就是说，不同群的这些 Hilbert 空间之间会互相「对话」！这称为 Langlands 函子性（Langlands functoriality）。本课程的部分目标，就是在物理中解释这个结构。

但 Langlands 纲领本身所做的，是取来这些自守形式，让它们在一块棱镜中做谱分解（TODO：画《月之暗面》封面上那块棱镜，或者——也许因为这事与物理有关——画《μ 子之暗面》？）。自守形式从左边进入，棱镜在 Hecke 代数（即所有这些两两交换的 Hecke 算子构成的代数）的作用下把它们谱分解。不同的「颜色」（特征值）由数域的 Galois 群的表示给出，这是一个出人意料又神奇的事实。此外还有对偶性：这些 Galois 表示并不是映入 $G$ 的复点，而是映入 $G^\vee_{\mathbb{C}}$，其中 $G^\vee$ 是在所谓 Langlands 对偶下与 $G$ 对偶的群。

例如，若 $G=\mathrm{PSL}_2(\mathbb{R})$，则 $G^\vee=\mathrm{SL}_2(\mathbb{C})$。要看它如何进入故事，一个相对具体的办法是：若 $E$ 是 $\mathbb{Q}$ 上的椭圆曲线，则 $H^1(E)$ 是一个二维向量空间，带有 $\mathrm{Gal}(\overline{\mathbb{Q}}/\mathbb{Q})$ 的作用，而这是模形式出现在故事中的一部分方式。这就是该理论中的一种「颜色」（特征值）。从精神上说，这个表示的像落在 $\mathrm{SL}_2(\mathbb{R})$ 中，不过要精确看清这个想法，需要把设定搭得更仔细一些。椭圆曲线在这里以两种方式出现：椭圆曲线的模空间本身就是一个算术局部对称空间。看清高层概览并不需要这一点，但要真正证明任何东西，它都至关重要！知道椭圆曲线有一些覆叠是自守曲线，会很有用，因为这为 Langlands 纲领的两侧搭起了一座桥。这很有用，但只适用于 $\mathrm{GL}_2$——一般情形下你没有这座桥，两侧相距极远。例如，$\mathrm{GL}_3(\mathbb{Z})\backslash\mathrm{GL}_3(\mathbb{R})/\mathrm{O}_3$ 不是任何东西的模空间：它不是流形。这使你的证明困难得多，

\srcpage{4}

而对偶也越发神秘：为什么这些空间上的函数论会与 Galois 表示有任何关系？\footnote{另一个不必操心的技术细节：当我们把 $\mathbb{C}$ 换成 $\mathbb{Q}_\ell$（要把其中一些东西说精确，这一步是必要的）时，必须使用 étale 上同调，而不是奇异上同调或 Zariski 上同调。但这对本讲的要点并不关键。}

本课程的主要目标之一，是展示（几何）Langlands 纲领如何在物理中出现——不是量子力学，而是四维（拓扑）场论。我们不是从一个量子力学系统出发，而是把它换成一个丰富得多、复杂得多——也吓人得多——的东西。关键的形容词「拓扑」可以缓解这种恐惧：我们把动力学扔掉，只看 Laplacian 的基态，就像只看调和形式而不看其余一切。我们将尝试把 Langlands 纲领的结构，与这个 TFT 的结构对应起来。

为什么是四维？量子力学似乎足够典范，四维物理则不然。我们引入另一个形容词：算术量子场论（arithmetic quantum field theory），沿用算术拓扑的范式。这个想法是在数域与物理中出现的几何对象（常常是流形）之间做类比。只要类比足够牢固，你就可以设想：对流形做的构造，在数域的世界里也有意义。这套理论的基本教义，Weil 的 Rosetta 石（TODO：引用）中已有概述；它建立了数域、函数域与 Riemann 曲面之间的字典。

\begin{itemize}
  \item 给定数域 $F/\mathbb{Q}$，我们考虑 $\Spec(\mathcal{O}_F)$，它的点以 $\mathcal{O}_F$ 中的素元来标记。
  \item 数域与函数域之间的类比更古老、也理解得更好。我们把 $F$ 换成有限域 $\mathbb{F}_q$ 上的（光滑射影）曲线 $C$。$C$ 上的有理函数域 $\mathbb{F}_q(C)$ 有许多让人想起 $F$ 的结构，正则函数环 $\mathbb{F}_q[C]$ 则像 $\mathcal{O}_F$（例如两者都是 Dedekind 整环）。$C$ 的点就像 $\mathcal{O}_F$ 中的素元。
  \item 但为什么要停在 $\mathbb{F}_q$？设 $\Sigma$ 是紧 Riemann 曲面。在 Weil 的字典中，$\Sigma$ 的点是素元的类比物，人们可以尝试把数论问题做成几何的类比。$\Sigma$ 上的亚纯函数域类比于 $F$ 与 $\mathbb{F}_q(C)$；整数环的类比物则稍嫌复杂——$\Sigma$ 上没有非常值整函数，所以我们必须挖掉一些点，它们类比于数域设定中的无穷远点。
\end{itemize}

算术拓扑类比的关键变化，是把 Riemann 曲面换成 3-流形。这一出人意料的变化，理由在于：Riemann 曲面与代数闭正特征域上的曲线有很强的相似性。当你研究一个点 $\Spec \mathbb{F}_q\hookrightarrow C$ 时，你应当记住 Galois 群作用所给出的内部结构。étale 拓扑告诉我们，要把 $\Spec(\mathbb{F}_q)$ 想成某种圆，因为 $\Spec(\mathbb{F}_q)$ 的 étale 基本群是 $\Gal(\overline{\mathbb{F}}_q/\mathbb{F}_q)\cong\widehat{\mathbb{Z}}$。Frobenius 是这个基本群的一个拓扑生成元。所以这里的结构太多，无法与一个 Riemann 曲面相匹配。若做基变换到 $\overline{\mathbb{F}}_q$，把这些「圆」换成它们的「万有覆叠」，我们就得到一个额外的线方向；拓扑/上同调看不见它，因为它是可缩的，但现在我们得到的东西，感觉上更像一个 Riemann 曲面了。

因此，从 Galois 理论的观点看，$\mathbb{F}_q$ 上曲线的函数域，感觉上不太像 Riemann 曲面，而更像 Riemann 曲面在 $S^1$ 上的丛。这与一个 Riemann 曲面 $\Sigma$ 连同微分同胚 $\varphi:\Sigma\to\Sigma$ 是等价的数据。然后我们把丛造出来：
\[
\Sigma\times[0,1]\big/\big((x,0)\sim(\varphi(x),1)\big),
\]
这个构造称为映射环面（mapping torus）。在函数域这一侧，我们把 Frobenius 想成 $\varphi$。这里有个差别：在 Riemann 曲面这一侧，我们并没有 $\varphi$ 的典范选择（恒等映射太无聊，还是别用它），所以我们把 $\varphi$ 想成「一般的」。

于是我们得到算术拓扑的字典，它由许多人共同建立，包括写信给 Mazur 的 Mumford，以及 Mazur、Kapranov、Reznakov、Morijita 和 Kim\yizhu{「Reznakov」「Morijita」疑为 Reznikov、Morishita 之误；此处照原文拼写。}{A（原文 p.4）}。这又称为「纽结与素数」字典：数域是 3-流形的类比物。这不是凭空冒出来的东西，而是 Weil 字典的一个精细化。

正如并非所有数域都有唯一的 Frobenius（该说 Frobenius 们吗？），而是在不同素数处有不同的 Frobenius，我们的 3-流形 $Y$ 也不只是曲线上的曲面丛。数域这边的素元，现在对应于 $Y$ 中的嵌入圆周，即纽结。局部域，例如 $\mathbb{Q}_p$，看起来像 2-流形。能纤维化到圆周上的 2-流形并不多，但这没关系。这个类比还有许多别的方面，

\srcpage{5}

例如 Legendre 符号与环绕数（linking number）之间的关系，等等。nLab 上关于算术拓扑的页面有一份很好的清单。

这并不是一个极其精确的字典，也不要犯这样的错误：把具体的素元与具体的纽结对应起来。例如，如果你说 $\mathbb{Q}$ 就是球面，那么你会发现 Poincaré 猜想在数域设定下不成立，那就不妙了。不如这样设想：数域是 3-流形的一类新例子，与通常的 3-流形既有共同点，也有别的性质；函数域则是另一族例子。于是我们就可以研究这类新的、丰富的例子。

回到为什么是四维拓扑场论的问题。首先我们得讨论拓扑场论究竟是什么；但我们将会看到，这种对象的基本不变量之一是：对每个 $(n-1)$-流形，都得到一个向量空间。所以 Langlands 纲领把向量空间（或与之相关的东西，例如分次向量空间、链复形）指派给函数域与数域；在我们的类比中它们是三维的，因此我们预期物理中有一个四维的故事。

更一般地，一个 $n$ 维量子场论是有动力学的：除了 $(n-1)$-流形上的向量空间，你还得到一个 Hilbert 空间结构与一个 Hamiltonian（哈密顿量）。同样，你手里的可能不是向量空间而是链复形之类的东西。但 Hamiltonian 使这件事更像是在你的余维 1 流形上的一个量子力学问题。在拓扑理论中，Hamiltonian 是 $0$。

现在回到局部对称空间：若 $F$ 是数域，我们认为场论指派给它的向量空间是 $L^2(\Gamma_{\mathcal{O}_F}\backslash G/K)$。空间 $\Gamma\backslash G/K$ 并不直接出现；相反，它是 3-流形上某些相关方程之解的模空间。

故事的其他部分也能搬过来。打开导子/分歧 $N$，我们就不只有一个 3-流形，还有其中的一个纽结或链环，我们把它想成奇点可能出现的位置。在物理中，这些称为「余维 2 缺陷」，是一般量子场论中的一项重要数据。

回顾一下：今天我们讲得很快，周四也会讲得很快。但从下周起，课程大多会放慢，那时我们会更仔细地追踪这个故事两侧的结构，努力在数论与物理这两种范式之间那条安全的几何钢丝上保持平衡。

周四我们会深入探讨数论这边可以摆弄的那些变量在物理中的类比物：改变导子会怎样，改变数域 $F$ 会怎样，玩一玩函子性又会怎样，等等。
